% common-scrbook/preamble/segment.tex
%
% Segment-level commands for the scrbook target. The kramdown converter
% (Mono::TypesetScrbook / AsfScrbookLatex) emits \segmenthead /
% \segmentappendixchapter at segment boundaries and \segmentfield at each
% segment-internal field heading (Formal Expression, Epistemic Status,
% Discussion, etc.). This file defines those commands as vanilla scrbook
% structures — no Tufte machinery, no tinted boxes, no margin notes.
%
% Status and stage arguments are present in the macro signatures (so the
% emitter doesn't need a separate scrbook code path) but ignored visually:
% scrbook's promise is "the legible read," and per-segment epistemic
% metadata lives in the chunk-format / assembled markdown for anyone who
% wants it. Don't decorate the page with discipline metadata.

%----------------------------------------------------------------------------
% \segmentrulewidth — referenced by AsfLatex's table converter when a
% narrow-area table needs to escape into the margin column. The kaobook
% target defines this as \textwidth + \marginparsep + \marginparwidth
% (body + margin gutter + margin column) so escaping tables visibly step
% into the Tufte margin. In scrbook there's no margin column to step
% into, so we collapse to \textwidth — escaping tables stay at body width
% and the per-column-weight heuristic in segment_renderer.rb still
% distributes columns sensibly across the body band.
%----------------------------------------------------------------------------
\newcommand{\segmentrulewidth}{\textwidth}

%----------------------------------------------------------------------------
% \segmenthead{Type}{Title}{Status}{Stage}
%
% In-part segment header. Renders as a vanilla \section, with the Type
% rendered as an italic genus marker preceding the title. Reads in the ToC
% and on the page as:
%
%     1.1  Definition: Agent-Environment Coupling
%
% with "Definition:" italicized — quiet but informative. Status and stage
% arguments are accepted but unused in this target.
%----------------------------------------------------------------------------
\NewDocumentCommand{\segmenthead}{m m m m}{%
  \section[\textit{#1:} #2]{\textit{#1:}~#2}%
}

%----------------------------------------------------------------------------
% \segmentappendixchapter{Type}{Title}{Status}{Stage}
%
% Appendix-segment header — appendix segments are chapter-level entities
% per the four-volume hierarchy decision. Same italic-prefix discipline
% as \segmenthead, just at chapter level. Numbering follows scrbook's
% native \appendix machinery (A, B, C, …).
%----------------------------------------------------------------------------
\NewDocumentCommand{\segmentappendixchapter}{m m m m}{%
  \chapter[\textit{#1:} #2]{\textit{#1:}~#2}%
}

%----------------------------------------------------------------------------
% \segmentfield{Title}
%
% Segment-internal field marker — Formal Expression / Epistemic Status /
% Discussion / Findings. NOT a numbered \subsection (would inflate the
% ToC and read as 1.1.1, 1.1.2, …, which is the wrong visual register
% for what's structurally a labeling cue). Renders as a small italic-bold
% standalone leader, with a thin rule's worth of breathing room above
% and below.
%----------------------------------------------------------------------------
\NewDocumentCommand{\segmentfield}{m}{%
  \par\medskip%
  \needspace{3\baselineskip}%
  \noindent\textit{\textbf{#1.}}\par%
  \nopagebreak[4]\nobreak\smallskip%
}

% Backward compatibility: AsfLatex (the base class shared with the kaobook
% converter) emits \segmentsubhead in segment-mode rendering. The volume-
% mode AsfScrbookLatex overrides convert_header so this isn't invoked from
% the assembled-markdown path, but defining it as an alias keeps any future
% segment-mode call from erroring out.
\let\segmentsubhead\segmentfield

%----------------------------------------------------------------------------
% \workingnotessubhead{Title}
%
% Sub-question heading inside the workingnotes environment. Same visual
% register as \segmentfield but expected to appear inside the working-notes
% box (defined in environments.tex). The workingnotes box already shifts
% the body to \small, so this leader is small + italic to match.
%----------------------------------------------------------------------------
\NewDocumentCommand{\workingnotessubhead}{m}{%
  \par\medskip%
  \needspace{2\baselineskip}%
  \noindent\textit{\textbf{#1.}}\par%
  \nopagebreak[4]\nobreak\smallskip%
}

%----------------------------------------------------------------------------
% \eqtag{Tag-content}
%
% Claim-atom label. Source markdown looks like:
%
%     *[Derived (sector-condition, from prior-slug)]*
%
%     $$ \alpha > \rho/R $$
%
% or the same tag above a prose claim (no display equation). These are
% annotations on the next claim, not headings — they belong in the outer
% margin. The converter emits \eqtag immediately before the claim
% (display equation, or the prose block) so \marginpar aligns with its
% first line.
%
% typearea under DIV=12 / letter / twoside leaves ~98pt of outer margin;
% \marginparwidth / \marginparsep are set in setup.tex to use most of it.
% \marginnote rather than \marginpar: see setup.tex.
%----------------------------------------------------------------------------
\NewDocumentCommand{\eqtag}{m}{%
  \marginnote{%
    \raggedright
    \footnotesize\itshape
    \color{wn-fg}%
    #1%
  }%
}

%----------------------------------------------------------------------------
% \externalref{Slug}
%
% Cross-reference fallback for slugs that don't have a \label in this
% volume — typically references into sibling volumes (e.g., AAT segments
% referenced from the ELI volume). The kaobook target's xr-hyper plumbing
% (Workstream B) will eventually resolve these to numbered cross-volume
% refs; until then, we render the slug itself in monospace so the reader
% sees a legible pointer rather than LaTeX's `??` undefined-ref marker.
%----------------------------------------------------------------------------
\NewDocumentCommand{\externalref}{m}{\texttt{\##1}}

%----------------------------------------------------------------------------
% \statusbadge{Status} / \stageglyph{Stage}
%
% Visual chrome from the kaobook target — explicitly no-op here. The
% emitter still calls them (so the IAL parser doesn't need a target-aware
% branch), but they produce nothing in scrbook.
%----------------------------------------------------------------------------
\NewDocumentCommand{\statusbadge}{m}{}
\NewDocumentCommand{\stageglyph}{m}{}
\NewDocumentCommand{\segmentstage}{m}{}
\NewDocumentCommand{\segmentfoot}{}{\par\medskip}

%----------------------------------------------------------------------------
% \missingsegment{Type}{Slug}{Claim}
%
% Placeholder for slugs listed in OUTLINE.md without a corresponding src/
% file. Renders as a real \section so the ToC reflects the gap, plus a
% visible body marker so a reader scanning the chapter sees what's
% missing rather than missing the gap entirely.
%----------------------------------------------------------------------------
\NewDocumentCommand{\missingsegment}{m m m}{%
  \section[\textit{#1:} \texttt{\##2} (missing)]{\textit{#1:} \texttt{\##2} \emph{(missing)}}%
  \label{seg:#2}%
  \par\smallskip%
  \noindent\emph{Segment \texttt{#2} not yet written.}\par%
  \IfStrEq{#3}{}{}{\smallskip\noindent\textit{Claim:} #3\par}%
  \medskip%
}

%----------------------------------------------------------------------------
% \missingappendixchapter{Type}{Slug}{Claim}
%
% Same pattern as \missingsegment but at chapter level — for missing
% appendix segments, since appendix segments are chapters in our
% hierarchy.
%----------------------------------------------------------------------------
\NewDocumentCommand{\missingappendixchapter}{m m m}{%
  \chapter[\textit{#1:} \texttt{\##2} (missing)]{\textit{#1:} \texttt{\##2} \emph{(missing)}}%
  \label{seg:#2}%
  \par\smallskip%
  \noindent\emph{Segment \texttt{#2} not yet written.}\par%
  \IfStrEq{#3}{}{}{\smallskip\noindent\textit{Claim:} #3\par}%
  \medskip%
}

%----------------------------------------------------------------------------
% \outlinegap{Description}
%
% Gap marker — for --GAP-- entries in OUTLINE tables (theory not yet
% developed; no slug, just a description of the open question). Renders
% as a flagged inline notice rather than a structural heading.
%----------------------------------------------------------------------------
\NewDocumentCommand{\outlinegap}{m}{%
  \par\medskip%
  \noindent\textbf{[Gap]}~\emph{#1}\par%
  \medskip%
}

%----------------------------------------------------------------------------
% Stub macros referenced by main.tex / kaobook converter that have no
% scrbook analog. Define as no-ops so the converter is target-agnostic.
%----------------------------------------------------------------------------
\NewDocumentCommand{\asfChapterFormat}{}{}
\NewDocumentCommand{\asfAppendixToCremap}{}{}
